The project set out to design and develop a hybrid chatbot for medical haematology using sequential models for intent classification and retrieval-based response generation. That aim was achieved. A substantial artefact was implemented that integrated BiLSTM classifiers, curated retrieval, safety rules, model switching, PostgreSQL logging, a browser UI, an admin monitoring console, a retraining workflow, and split-aware evaluation.

From a research perspective, the project demonstrated that sequential models remained appropriate for bounded intent-classification tasks when the system scope was clear and the surrounding data and retrieval infrastructure was strong. From an engineering perspective, the project demonstrated that a small-scale MLOps implementation could still include important lifecycle concepts such as data staging, version-aware artefacts, train-validation-test separation, automated testing, monitoring, and feedback-based retraining.

The final results were credible rather than inflated. Test accuracy remained close to or above 0.90 across both model profiles, macro F1 remained above 0.86, and the full automated test suite passed with 52 tests. At the same time, the project acknowledged its constraints: moderate dataset size, English-only operation, rule-based entity refinement, and the deliberate exclusion of diagnosis and unrestricted patient-specific interpretation.

Several directions for further work were identified. Additional real user phrasing should be collected to reduce reliance on synthetic paraphrases and to improve communication-class robustness. More report-oriented entity coverage could be added, especially for fine-grained blood count flags, morphology terms, unit-aware parsing, and richer multi-parameter report interpretation. Formal usability testing with laboratory users should be conducted using structured instruments. The retrieval layer could be expanded with richer ranking and explanation support. Finally, multilingual support and document-grounded report assistance could be investigated, provided that the safety boundary remained explicit.

Overall, the project satisfied the goals of a final-year software and machine learning dissertation. It combined theoretical grounding, technical implementation, operational evaluation, and critical reflection into one coherent body of work.
